\section{The sixteen remaining chain models}\label{app:cases}

We record the models, exceptional parameters and completeness arguments for
Types~1--9 and 11--17. The notation $(m,\mathcal A,\mathcal B)$ means that
$A$ has the central chain of length $2m$, and its other chains occur in
reflected pairs with lengths in $\mathcal A$; the reflected chain pairs
of $B$ have lengths in $\mathcal B$. All shifts have length one. A lower-case
center parameter denotes the square of a rational center. We use the moment polynomials from \eqref{eq:chains}.
For three singleton squares whose first power sums are $p_1,p_2,p_3$,
we always use the exact monic polynomial
\[
 K(X)=X^3-p_1X^2+\frac{p_1^2-p_2}{2}X
       -\frac{p_1^3-3p_1p_2+2p_3}{6}.
\]
Thus statements that this polynomial splits into squares include the
individual center conditions. A finite-field test throughout this appendix
retains repeated roots and every possible pole chart unless a stated
valuation argument excludes that chart. An overlap, a zero reflected center,
or a repeated singleton invalidates the proposed input.

For Types~1, 9, 14 and 17, let $a=u^2$ be the noncentral $A$ center
square, $b=v^2$ the $B$ doubleton center square, and set
\[
 p_k=C_m(k)+M_{5-m}(k,a)-M_2(k,b).
\]
The necessary noncollision branches for $m=1,2,3,4$ are respectively
\[
 a=5b+4,\quad a=(10b+9)/4,\quad a=(5b+3)/3,\quad a=(5b+1)/4.
\]
Every other factor of the fourth Newton residual forces an overlap.
The factors $b=0$ give coincident reflected chains. The polynomial identities,
including every discarded factor, are checked in
\cert{check_shift_chains.py} and \cert{check_all_shift_chain_types.py}.

\subsection{Type 1}\label{app:type1}
\noindent\emph{Arithmetic details in Section~\ref{detail:type1}.}\enspace
Here $(m,\mathcal A,\mathcal B)=(1,(4),(1,1,1,2))$ and $a=5b+4$.
The three singleton squares are the roots of $K$ above. Their product gives
a point on an elliptic curve with minimal coefficients
$[1,0,0,-51865,3207225]$. Its rational group is torsion-free of rank four.
The reductions of the full group modulo doubling at $7,67,73,107,257$,
with both integral and pole tests, leave two classes represented by $b_0=0,1$.
At 73 these classes require $b\equiv b_0$. This is a restriction to global
parity classes, rather than a bound on the height of $b$.

The root-field character turns each class into a cover over
$K=\mathbf Q(w)$, where $w^3-w^2-330w-149578=0$ and
$r=(-w^2+20w-5426)/7168$. For $d=5(b_0-r)$ the exact model is
\[
 E_d:Y^2=X(X-dr)(X+4d/5),\qquad X=db.
\]
It has full two-torsion. A split two-descent and two nondegenerate
four-dimensional Cassels pairing minors give rank bounds two and one for
$b_0=0,1$, respectively. In the second case the seed is non-torsion.
For the three embeddings at 73 the primitive logarithm and parameter tangent
are
\[
 L=(30,34,22),\qquad V=(55,61,36),\qquad
 (L_iV_j-L_jV_i)_{i<j}=(33,16,28).
\]
At least two of the three signs agree. The corresponding logarithmic
linear combination has a unit derivative, so its only zero in the entire
disk is $b=1$. The doubleton then shares a vertex $1/2$ with the central
chain.

For $b_0=0$ write $b=v^2$. The center square $5b+4$ and the local character
force $v_2(v)\geq3$ if $v\ne0$. Use the sextic
\[
 h(B)=-800B^6+2625B^5-12050B^4+529B^3+54448B^2+53744B+14400.
\]
The singleton discriminant gives $h(b)$ square. In $F=\mathbf Q(\theta)$,
$h(\theta)=0$, the class of $1-b/\theta$ belongs to the norm-square
$S$-unit space with
$S=\{2,3,5,23,41,2111,11801,144461\}$. The local conditions at these
primes, at 73 and at the real places force this class to be trivial.
A second division at the two dyadic places then excludes the primitive
logarithm direction forced by every nonzero $v$. The computation uses a
certified class group of discriminant $14755078257381100$, complete local
squareclass bases and a four-cover linear system. Thus the only parameter
in this disk is $v=0$, which repeats reflected chains. The odd-prime
linear term estimates and the dyadic primitive direction estimates apply
to the whole corresponding disk.
The group, parity, pairing and local calculations are given in
\cert{check_type1_parity.m}, \cert{type1_descent/verify.gp},
\cert{type1_descent/tangents.gp}, \cert{type1_descent/disk73.m},
\cert{type1_descent/sextic.gp}, \cert{type1_descent/sextic_logs.m} and
\cert{type1_descent/sextic_four.gp}.

\subsection{Type 2}\label{app:type2}
\noindent\emph{Arithmetic details in Section~\ref{detail:type2}.}\enspace
The type is $(1,(1,3),(1,1,3))$. If $u,v$ are the length-three
centers on the two sides, the noncollision factor of the moment equations is
\[
 8u^4-40u^2v^2-58u^2+80v^4+50v^2+95=0.
\]
Put $y=8u^2-20v^2-29$. This gives
\[
 y^2=-240v^4+760v^2+81.
\]
For $v\ne0$, the substitution
$X=9(y+9)/(8v^2)+95/2$, $Y=9X/(2v)$ identifies its smooth projective
model with $E:Y^2=X^3-95X^2+2560X$. The two exceptional points
$(0,\pm9)$ give the identity and $(0,0)$. The two-isogenous curve is
$Y^2=X^3+190X^2-1215X$. For an isogeny descent class $d$, a point
requires a primitive integral solution of
\[
 w^2=du^4+Au^2v^2+(B/d)v^4.
\]
On $E$ the only possible classes are $1,2,5,10$ by positivity; 2 and 5
are excluded modulo 3. On the dual curve, among $\pm1,\pm3,\pm5,\pm15$,
all except 1 and $-15$ are excluded in the primitive charts modulo
16 or 128. Both descent images therefore have order two and the rank is
zero. Good reduction at 7 and 11 bounds torsion by ten, and the ten points
\[
 O,(0,0),(16,\pm144),(40,\pm120),(64,\pm192),(160,\pm1440)
\]
attain this bound. The complete quartic list is
$(v,y)=(0,\pm9),(\pm1/2,\pm16),(\pm3/2,\pm24)$.
At $v=0$ the value of $u^2$ is $19/4$ or $5/2$; neither is square.
At the other values the $B$ tripleton meets the central $A$ chain.
There is no rational point at infinity since $-240$ is nonsquare.
The elementary residue tables and every substitution are certified by
\cert{check_type2_exclusion.py}; \cert{check_type2_exclusion.m} gives
an additional arithmetic check.

\subsection{Type 3}\label{app:type3}
\noindent\emph{Arithmetic details in Section~\ref{app:type3-details}.}\enspace
The type is $(1,(1,3),(1,2,2))$. Let $a,b$ be the singleton and
tripleton center squares on $A$, $c$ the $B$ singleton square, and
$S,T$ the sum and product of its doubleton squares. Moment elimination gives
\[
 1024b^2-1520bc+636b+640c^2-1860c+2025=0,
 \qquad S=(20b-8c+27)/6.
\]
To justify division in the second equality, the exceptional resultant is
\[
 (4c-25)(4c-9)(4c-1)^4
 (80c^2-280c+281)(144c^2-616c+3025)(144c^2-328c+4489).
\]
The three quadratics have negative discriminant. The constant possibilities
$c=1/4,9/4,25/4$ force overlaps or nonsquare $b$.
Set $w=(b-25/4)/(c-25/4)$ and $z=(w-16)/(w-1)$.
The direction $w=1$ gives $T=0$, and $c=25/4$ was already treated.
With
\[
\begin{aligned}
 L&=z^2-57z+1656,& A&=49z^2-393z+39744,\\
 B&=9z^2+487z+14904,& C&=9z^2+247z+39744,\\
 V&=2z^2-14z+1587,
\end{aligned}
\]
one has $(a,b,c)=(A,B,C)/(4L)$ and $T=300zV/L^2$.
Positivity gives $z>0$. The individual square conditions and the
resultant support bounds reduce the squarefree parts of $L$ and the
numerator times denominator of $z$ to
$(d,k)=(1,6),(115,46)$. The projective tests use $2^{13}$ and $3^7$,
and the support bounds $d\mid690$, $k\mid138$ come from resultants,
so no other global squareclass can occur.

The first pair gives a point on
$E_6:Y^2=X(X^2-342X+59616)$ with $X=6z>0$.
This curve has rank zero and torsion of order two, so its only finite
rational abscissa is zero. For the second pair put $t=1656/z$ and
$q(t)=t^2-131t/8+3381$. Then $t$, $10L(t)$ and $15q(t)$ are squares.
The rational rank-one curve $E_{10}$ fixes the class of $t-r$, where
$r=(57+15\sqrt{-15})/2$. Over $\mathbf Q(\sqrt{-15})$ use
\[
 \ell=64(16-r)\,3/5,\qquad
 E:Y^2=X\bigl(X^2+\ell q'(r)X+\ell^2q(r)\bigr),\qquad t=r+X/\ell.
\]
The point with $X=7200\sqrt{-15}-26400$, $Y=360X$, and $(0,0)$
generate a subgroup of finite odd index. This follows from rank one,
the complete two-torsion and nondivisibility of both translates at 19.
Elliptic Chabauty has index obstruction 4 and gives precisely the rational
images $t=16,\infty$. Back-substitution gives only overlapping chains.
The exact elimination and support bounds are in
\cert{check_type3_exclusion.py}; ranks, indices and rational images are in
\cert{check_type3_exclusion.m}.

\subsection{Type 4}\label{app:type4}
\noindent\emph{Arithmetic details in Section~\ref{detail:type4}.}\enspace
The type is $(1,(2,2),(1,1,3))$. Let $c$ be the $B$ tripleton square,
$S,T$ the symmetric functions of the $A$ doubleton squares, and $U,V$
those of the $B$ singleton squares. The unique valid branch is
\[
 S=(20c+25)/6,\quad T=(240c^2+40c+167)/48,
\]\[
 U=(44c+91)/12,\quad V=(256c^2+97c+466)/36.
\]
The exceptional factors $c=1/4,9/4$, the linear branch $S=2c+1/2$
and the quadratic branch in $S$ all give explicit common vertices.
The remaining possible loss of degree has $c=-11/4$, which is not square.
The original centers require $c$ and the three polynomials
\[
 q_T=3(240c^2+40c+167),\quad q_A=-80c^2+220c+31,
 \quad q_B=3(-720c^2+2152c+275)
\]
to be squares. Descent on $Y^2=q_Tq_Aq_B$ gives three fake classes.
The three individual quadratic component norms select the class of
$c_0-\theta$, $c_0=1/4$. In an identity
$c-\theta=\lambda(c_0-\theta)z^2$, unit valuations away from 2,3,5,
the split Newton polygons at 3 and 5, and positivity at the two real
roots of $q_A$ imply $\lambda\in\mathbf Q^{\times2}$.
In particular the common scalar in a fake class is removed using the
original individual squares.

Take $s^2-s-1=0$, $r=(14+27s)/20$ and $\gamma=1/4-r$.
Every valid configuration therefore lies over
\[
 z^2=\gamma(c-r)(720c^2+120c+501).
\]
Put $L=720\gamma$, $u=Lc$, $v=Lz$ to obtain its monic elliptic model.
It has rank one over $\mathbf Q(\sqrt5)$, a non-torsion point above
$c=1/4$, and torsion of order two. Saturating at two and adjoining
torsion gives an odd-index subgroup. Elliptic Chabauty returns index
obstruction 8 and complete images $c=1/4,\infty$. The finite image
makes the paired-center squares 1 and 4 and repeats central vertices.
An original center is finite. All auxiliary proof bounds are Minkowski
bounds. The complete calculation is \cert{check_type4_exclusion.m},
with independent identities and scalar normalization in
\cert{check_type4_scalar_normalization.py}.

\subsection{Type 5}\label{app:type5}
\noindent\emph{Arithmetic details in Section~\ref{detail:type5}.}\enspace
The type is $(1,(2,2),(1,2,2))$. For the singleton square $c$ of $B$,
let $S,T$ and $U,V$ describe the two doubleton squares on $A$ and $B$.
The noncollision branch is
\[
 S=(20c-85)/16,\quad T=(80c^2-840c+589)/256,
\]\[
 U=(12c-83)/16,\quad V=(4c-49)(4c-9)/256.
\]
The exceptional values $c=1/4,9/4,25/4$ give central or cross-color
collisions. The necessary individual squares are $c$ and
\[
 q_1=80c^2-840c+589,\quad q_2=80c^2-40c+4869,
 \quad q_3=80c^2-1064c+5125.
\]
Descent on $Y^2=q_1q_2q_3$ gives the two classes represented by
$c_0=49/4,81/4$. The normalization lemma proves that
$c-\theta=\lambda(c_0-\theta)z^2$ implies
$\lambda\in\mathbf Q^{\times2}$. For odd primes outside 5 choose a
quadratic factor with unit discriminant. Square norms then give even
valuations. At 5, $c$ is integral and the root of $q_3$ of valuation
$-1$ gives a square unit ratio. At 2, the conditions force $v_2(c)=-2$;
the split roots of $q_1$ make corresponding root-difference valuations
odd for both $c$ and $c_0$, so their ratios have even valuation.
Finally $c\geq49/4$ and the two real roots of $q_1$ lie below this,
which fixes the positive sign of $\lambda$.

For $c_0=49/4$, take $K=\mathbf Q(r)$, $q_1(r)=0$, and
$g(c)=(c-r)q_3(c)/80$, $\gamma=g(49/4)$.
The cover $z^2=\gamma g(c)$ is elliptic of rank one and torsion order
two over $K=\mathbf Q(\sqrt{505})$. An odd-index subgroup obtained
by two-saturation gives complete rational images $49/4,\infty$,
with index obstruction 16. At $49/4$, $V=0$, so a reflected doubleton
pair repeats itself. For $c_0=81/4$, let $q_2(r)=0$ and
$\gamma=c_0-r$. Use the extra square $c$ to obtain
\[
 z^2=\gamma c(c-r)q_3(c)
\]
over $K=\mathbf Q(\sqrt{-95})$. Inversion in $c$ gives a rank-one
elliptic curve, again with torsion order two. The complete rational
images are $0,81/4$ and the index obstruction is 8. The value zero
has $S<0$; at $81/4$ both inputs are the ten half-integers from
$-9/2$ to $9/2$. The cover has no additional rational infinity
for the original sextic, whose leading coefficient is nonsquare.
The certificates are \cert{check_type5_scalar_normalization.py} and
\cert{check_type5_exclusion.m}. These checks retain the scalar in the
fake descent until its squarehood has been proved.

\subsection{Type 6}\label{app:type6}
\noindent\emph{Arithmetic details in Section~\ref{detail:type6}.}\enspace
The type is $(1,(1,1,2),(1,4))$. Let $a$ be the $A$ doubleton square,
$S,T$ its singleton symmetric functions, $b$ the $B$ singleton square
and $c$ its quadrupleton square. Moment elimination gives the smooth cubic
\[
 H_6=a^3-5a^2c-9a^2+10ac^2+15ac+24a-10c^3+10c^2-20c-16=0,
\]
and
\[
 b=(5a-c+5)/4,\quad S=(-3a+15c+22)/4,
 \quad T=(8a^2-40ac-67a+80c^2+95c+149)/16.
\]
The eliminated factors $a-c$ and $(a-c-1)^2-4c$ are overlaps.
The possible vanishing coefficient
$E_6=3a^2-8ac-7a+6c^2+2c+4$ has resultant proportional to
$c(c-4)^2(c-1)^3$; these values repeat reflected chains or meet the
central chain. Thus the formulas apply to every valid input.
The square cover $c=v^2$, resolved at $(a,c)=(4,0)$ by $a=4+vw$,
has plane-quartic equation
\[
 vw^3-5v^2w^2+10v^3w-10v^4+3w^2-25vw+50v^2-40=0.
\]
The origin $(1,0)$ on $H_6$ gives minimal model
$y^2+xy+y=x^3+x^2-2913x+80031$, of rank two and torsion order two.
Exact division at 2 and 3 and simultaneous reduction modulo 216 leave
exactly the classes represented by $(1,1),(9,1),(16,4)$.
The pole charts and the extra branch condition $3w^2=40$ over $(4,0)$
are retained in this reduction.

For a common Prym construction put
\[
 t=(a-1)/c,\quad A=t^3-5t^2+10t-10,\quad B=-6t^2+5t+20,
 \quad Y=2Ac+B.
\]
Then $Y^2=5t(32t^2-155t+192)$. In $K=\mathbf Q(r)$,
$r^2=-551$, let $e=(155+r)/64$ and $h=-B(e)/(2A(e))$.
The character of $q_1=a-1-ec$ is fixed in each class modulo 216,
with representatives $-e,8-e,-e$. Adjoining its square root gives
an unramified double cover of the quartic and a genus-two Prym.
On one twist two independent collision divisors span the global
logarithm space; their tangent determinants isolate the classes of
$(1,1)$ and $(16,4)$.

For the remaining twist, affine descent on divisor classes of degree one
is needed. Its global affine space has dimension 14 after homogenization
with parity, the even subspace has dimension 13, and scalars account
for 11 dimensions. Two explicit classes $D,T$ span the quotient.
The verifier constructs complete local images, including dimensions
$8,6$ at the two dyadic places, and proves the requisite divisor
representability via local Brauer obstructions. On the disk of
$(9,1)$ the displacement, normalized by $(v-1)/8$, has primitive
row $(1,0,0,0)$ outside the global span. Opposite local sheets add
$(0,0,1,0)$, also outside that span. The integrated power-series
bounds apply to the whole disk, not just its first coefficients.
The only remaining value is $(a,c)=(9,1)$, a collision.
Rank, class and exact-division data are in \cert{check_type6_classes.m}.
The two Prym twists and their complete affine descent are in
\cert{type6_descent/class_bounds.gp}, \cert{type6_descent/prym.m},
\cert{type6_descent/beta_selmer.m} and \cert{type6_descent/beta_logs.m}.

\subsection{Type 7}\label{app:type7}
\noindent\emph{Arithmetic details in Section~\ref{detail:type7}.}\enspace
The type is $(1,(1,1,2),(2,3))$. Let $a,S,T$ describe the doubleton
and singleton squares on $A$, and $b,c$ the doubleton and tripleton
squares on $B$. Every valid configuration lies on
\[
\begin{split}
 H_7={}&256a^3-1280a^2c-64a^2+2560ac^2-2560ac-2976a\\
       &-960c^3+3920c^2+780c+3159=0,
\end{split}
\]
with
\[
 b=(10a-4c+15)/6,\quad S=(16a+20c+81)/12,
\]\[
 T=-(128a^2-640ac+48a+80c^2-1380c-1701)/144.
\]
The discarded factors are $4c-1,4c-9$, the overlap quadratic
$16a^2-32ac-8a+16c^2-8c+1$, and $8a-12c-5$.
On the last branch the next equation forces $c=-11/4$ unless there
is already a collision. The subresultant coefficient can also vanish
only at the two collision values, a negative square, or roots of
quadratics with negative discriminant. This treats the exceptional loci.

With origin $(1,1/4)$ the cubic has rank two and no torsion.
Saturation at 2 and 3 followed by reductions modulo 216 leaves
\[
 (1,1/4),\ (9,9/4),\ (4,49/4),\ (4,1/4),\ (1,25/4).
\]
The square cover $b=v^2$ has genus four and the rational branch point
$(1,25/4)$. Put $u=a-1$, $t=b/u$ and
$c=(25+(10-6t)u)/4$. After removal of the branch factor its equation is
\[
 (45t^3-145t^2+135t-27)u^2-8(45t^2-80t+21)u+144(5t-2)=0,
\]
so $y^2=5F(t)$ with $F=27t^3-79t^2+57t-9$.
The omitted origin is the pole $t=\infty$; the other point at $a=1$
has $t=2/5$. The five classes give respectively
\[
 (t_0,d)=(\infty,15),(2,2),(1/3,3),(3,3),(2/5,10),\qquad t=dz^2.
\]
The support bound $d\mid30$ follows from the valuations at zeros and poles.
For a root $e$ of $F$ over its cubic field, let
$F(t)=27(t-e)(t^2+g_1t+g_0)$ and $k=135d(t_0-e)$, with the
origin character used in the first row. The elliptic curves
\[
 Y^2=X(X^2+kg_1X+k^2g_0),\qquad X=kt
\]
have ranks $0,3,3,2,2$. The first is exhausted by its two-torsion.
For the rank-two rows at 61 and 389, the determinants of two global
logarithms and the parameter tangent are units, with residues 31 and
216. A second descent in the sextic algebra treats the $d=3$ row
at $t_0=1/3$ by a rank-one cover and transversality at 83.

For the final row $d=2,t_0=2$, the same descent produces an elliptic
curve over the sextic field $L=\mathbf Q(h)$,
$h^6-96h^4+144h^2-54=0$. Its model has rational two-torsion and a
two-isogenous curve. The paired isogeny Selmer spaces have dimensions
4 and 1. Locally at the unique dyadic place, the kernel on the
four-dimensional space has one nonzero class $s_*$. A conic lift of
$s_*$ supplies the second descent character. Its two possible local
halves both lie outside the global image. The power-series estimate
on $z=1+4\xi$ bounds every higher coefficient, and division by the
maximal local power of two forces one of those excluded halves unless
$\xi=0$. The resulting point $(9,9/4)$ is a collision.
All the other anchors give repeated or cross-color vertices.
The initial group and branch checks are \cert{check_type7_classes.m};
the complete cover, logarithm and divisibility data are
\cert{type7_descent/classes.m}, \cert{type7_descent/covers.m},
\cert{type7_descent/odd_logs.gp}, \cert{type7_descent/sextic.gp} and
\cert{type7_descent/dyadic.m}.

\subsection{Type 8}\label{app:type8}
\noindent\emph{Arithmetic details in Section~\ref{detail:type8}.}\enspace
The type is $(1,(1,1,1,1),(5))$. Scale the squared $B$ center and
singleton squares by four, writing $B=4b$. Define
\[
 P_B(X)=(X-B)((X-B-4)^2-16B)((X-B-16)^2-64B),
 \qquad F(B,X)=\frac{P_B(X)-P_B(1)}{X-1}.
\]
Every valid input requires the four roots of this monic quartic to be
rational squares. Its constant term is $5f(B)$, where
$f(T)=T^4-18T^3+114T^2-97T+405$. Denominator valuations imply
that odd denominator primes are 1 modulo 20 and that $B$ is integral
at two. Indeed the denominator reduction has roots $1-\zeta^i$;
splitting forces fifth roots of unity, and square ratios force $-1$
to be square. The exceptional prime 5 is ruled out by the constant
coefficient valuation, and at two the denominator polynomial has no root.
The genus-two two-cover descent fixes one of the classes represented by
$B_0=1,9,25$.

At 23, the class $B_0=1$ forces $B\equiv1$, and
$F(1,X)=(X-1)(X-9)^2(X-25)$. Write $U=B-1,V=X-9$.
Its quadratic tangent cone is $-64(2V^2-15U^2)$. Since 30 is
nonsquare modulo 23, its valuation at a nonzero pair is exactly twice
the minimum valuation, while higher terms have larger valuation.
Thus the entire singular disk forces $B=1$, a repeated singleton.

For $B_0=9,25$, let $r$ be a root of $f$, $q(T)=f(T)/(T-r)$,
$g(T)=T^3q(1/T)/q(0)$, $t=1/B$, $t_0=1/B_0$, $k=g(t_0)$.
The reciprocal cover has monic model
\[
 y^2=x^3+kg_2x^2+k^2g_1x+k^3g_0,\qquad x=kt.
\]
Use $K=\mathbf Q(h)$, $h^4-h^3+28h^2-54h+231=0$, with
$r=(-2h^3+h^2-33h+129)/15$. The ordered-pair root field has degree
12, signature $(0,6)$, discriminant $5528109^5$ and class group of
order three. The explicit formula, with support $2^{22}$ and bound
$2^{17}$, has margin greater than $0.862$; exact ideal reduction lowers
the bound to 1000. Both reciprocal Selmer groups have dimension three.
Three explicit point logarithms establish the matching rank three.

For $B_0=9$ the primes $17,131,197,269$ fix the center residue.
Two- and three-saturation give all global coefficient classes modulo
16 and 9. Their simultaneous reductions make the three local signs
at 131 equal. The three-row global logarithm matrix has a unit
transverse tangent determinant. For $B_0=25$ the unique unramified
quartic dyadic place suffices, and its odd-square parameter lies in
one disk modulo eight. Here the common valuation of the three-row
minors is one, and adjoining the rational tangent does not raise it.
In each case the resulting annihilating form has a unit linear term;
the integrated binomial expansion makes every nonlinear coefficient
strictly smaller throughout the disk. Hence $B=9$ or 25. At both
values $P_B(1)=0$, so the two sets share a vertex square.
The certificates are \cert{check_type8_product.m},
\cert{type8_descent/node.m}, \cert{type8_descent/factor_base.gp},
\cert{type8_descent/covers.m}, \cert{type8_descent/logarithms.m} and
\cert{type8_descent/signs.m}.

\subsection{Type 9}\label{app:type9}
\noindent\emph{Arithmetic details in Section~\ref{app:type9-details}.}\enspace
Here $(m,\mathcal A,\mathcal B)=(2,(3),(1,1,1,2))$,
$a=(10b+9)/4$. In homogeneous coordinates $[R:B:W]$, the singleton
cubic is
\[
\begin{split}
 64R^3-(352B+688W)R^2+(432B^2+1840BW+2124W^2)R\\
 \qquad+216B^3-492B^2W-1974BW^2-2025W^3=0.
\end{split}
\]
The matrix
\[
 \begin{pmatrix}-44&-84&347\\-1088&3102&4064\\16&6&-28\end{pmatrix}
\]
carries it to $E:y^2+y=x^3+x^2+295x+1109$. The full rational
group is generated by $G=(31,202)$ and has no torsion.
Three roots in a common finite fiber give collinear points $n_iG$,
so $n_1+n_2+n_3=0$. Put $M=\max|n_i|$. One index has absolute
value at most $M/2$. Explicit denominator and duplication bounds
compare the three naive heights in a common fiber and give
$-46\leq h(2P)-4h(P)\leq17$. Telescoping gives bounds for the
canonical height, and the proportional index gap gives $M\leq25$.
The resulting finite group-index list is therefore an exhaustive
consequence of a global height argument. Exact substitution in the
pencil leaves only $b=4$, or $b=0$ when repeated roots are allowed.
The first is the coincident half-integer configuration; the second
has a zero reflected center. The group computation is
\cert{check_type9_group.m}; the height constants, duplication identities
and complete fiber test are \cert{check_type9_exclusion.py}.

\subsection{Type 11}\label{app:type11}
\noindent\emph{Arithmetic details in Section~\ref{app:type11-details}.}\enspace
The type is $(2,(1,2),(1,2,2))$. Let $a,b$ be the $A$ singleton
and doubleton squares, $c$ the $B$ singleton square and $S,T$ the
symmetric functions of the two $B$ doubleton squares. The necessary
noncollision conic is
\[
 16b^2+80bc-140b-80c^2+280c-281=0,
 \qquad S=(5b-4c+16)/3.
\]
The second formula can fail only at $c=1/4,9/4,49/4$.
The first two values meet the central chain. At $49/4$ the only square
choice is $b=9$, and the two residual possibilities for $S$ repeat
central $A$ vertices. Put $w=(b-9)/(c-49/4)=m/n$ in lowest terms.
The square conditions force
\[
 m^2+5mn-5n^2=-5k^2,\qquad m:n\equiv0:1\text{ or }4:1\pmod{37}.
\]
To obtain this conclusion the resultant support restricts the squarefree
part to $-1$ or $-5$; the two residues for the $-1$ class fail the
remaining product-square test at 37. The denominator is prime to 37.

Use $r=(47w-40)/(4w-5)$. The excluded value $4w=5$ would give $T=0$.
The genus-three curve needed for the last class is
\[
 C:Y^2=r(r^2-11r-1)(7r^2-2r+18)(9r^2-69r-49).
\]
At 37 only the disks of $(0,0),(8,\pm1500)$ are possible.
The Jacobian has unconditional rank at most two. The two rational
classes
\[
 P=[(8,1500)-\infty],\quad
 D=\left(r^3-\tfrac{51}{4}r^2+\tfrac{23}{4}r+8,
          \tfrac{415}{8}r^2+\tfrac{785}{8}r+20\right)
\]
have independent logarithms, since
\[
 \log(58624P)/37\equiv(25,28,33),\qquad
 \log(58624D)/37\equiv(3,30,12)\pmod{37}.
\]
Their cross product is $(12,21,0)$. The corresponding annihilating
differential has unit constant coefficient in all three disks, with
values 32 at $r=8$ and $12/31$ in the $Y$-parameter at zero.
Its integral is injective on each whole disk by the bound
$k-1-v_{37}(k)>0$ for every nonlinear degree $k\geq2$.
Thus the three rational anchors exhaust the eligible disks.
At zero $T=0$; at $r=8$ one has $c=1/4$, a central collision.
The exact rank bounds and local arithmetic are
\cert{check_type11_genus3_rank_bounds.m} and
\cert{check_type11_thirty_seven_adic.py}; the global squareclass reduction
is included in \cert{check_all_shift_chain_types.py}.

\subsection{Type 12}\label{app:type12}
\noindent\emph{Arithmetic details in Section~\ref{detail:type12}.}\enspace
The type is $(2,(1,1,1),(1,4))$. With $b$ the quadrupleton square
and $c$ the singleton square on $B$, the fourth Newton residual is
$-4b(b-4)(b-1)(b+4c-10)$. The first three factors repeat reflected
chains or meet the central chain, so $c=(10-b)/4$. The singleton
polynomial is defined exactly by
$p_k=M_4(k,b)+((10-b)/4)^k-C_2(k)$, $k=1,2,3$.
Its product and square conditions reduce the global parity classes to
$b_0=1,9$, with $(b-r)/(b_0-r)$ square in the cubic field
\[
 K=\mathbf Q(w),\quad w^3+38w-504=0,\qquad r=(4-w)/8.
\]
Put $q(B)=B^2+(r-3/2)B+r^2-3r/2+43/32$.
The product polynomial is $f(B)=160(B-r)q(B)$.
For $b_0=1$ take $d=-160(1-r)$ and
\[
 E^{(1)}:y^2=x\bigl(x^2+dq'(10)x+d^2q(10)\bigr),\qquad x=d(b-10).
\]
For $b_0=9$ let $D=160(9-r)$, $c_0=10Dq(0)$, and
\[
 F(U)=(U-c_0/10)(U^2+10Dq_1U+10Dc_0).
\]
The model is $y^2=F(x+c_0/10)$ with $x=c_0/b-c_0/10$.
Since $b$ and $10-b$ are squares, $b\ne0$ and no genuine point is
lost in inversion. Both curves have rank one over $K$. Their Selmer
groups have order 16 and rational two-torsion order two, and the
Cassels--Tate pairing on the two-dimensional quotient by the seed
and torsion is nonzero. The first descent uses a certified sextic
field bound of 100 after field identification, and the second uses
Minkowski bounds throughout.

At 907 no pole is possible, since $-1$ is nonsquare. The three
embeddings and the full singleton splitting tests leave only
$b\equiv b_0$. The logarithm/tangent minors for the three pairs
of embeddings are respectively $(698,516,760)$ and $(238,499,300)$.
All are nonzero modulo 907. Two signs among three agree, so one
minor supplies a function with a unit linear coefficient and integral
derivative. Integration bounds prove its unique zero in the whole
disk is $b=b_0$. Both values are collisions.
The certificates are \cert{check_type12_center_parity.m},
\cert{check_type12_fields.gp} and \cert{check_type12_exclusion.m}.

\subsection{Type 13}\label{app:type13}
\noindent\emph{Arithmetic details in Section~\ref{app:type13-details}.}\enspace
The type is $(2,(1,1,1),(2,3))$. The $B$ doubleton and tripleton
squares $b,c$ satisfy $6b+4c=25$ after removal of the collision
factors $(b-1)(4c-9)(4c-1)$. The three singleton squares on $A$
have polynomial $G_c$ defined by
$p_k=M_2(k,(25-4c)/6)+M_3(k,c)-C_2(k)$.
The product elliptic curve has two global Kummer classes, denoted
$\alpha$ and $\beta$, with $\alpha=25/4$ and $\beta=1/4$.
The cubic field may be written
\[
 K=\mathbf Q(\theta),\qquad
 \theta^3-\tfrac{20}{3}\theta^2+\tfrac{13}{12}\theta-\tfrac{95}{48}=0.
\]
For the $\alpha$ class the two-isogenous model is
$Y^2=X(X^2+aX+e^2)$, where
$\delta=\alpha-\theta$, $d=-6\delta$, $e=d\alpha$ and
$a=2d(\alpha-2\theta)$. Its center map is
\[
 c=\frac{\alpha(X^2+aX+e^2)}{(X-e)^2},\qquad
 \frac{\alpha}{c}=1+24\delta^2t^2,\quad t=-X/Y.
\]
A valid lift has $XY(X-e)\ne0$. Over the unique ramified cubic
dyadic completion, $\delta t=\pm v/(4w)$ where $c=w^2$ and
$4w^2+6v^2=25$. The local divisibility character is injective on
the global Selmer image. Repeated halving to the primitive point
then requires one of two local halves that are both outside that
image. This excludes the $\alpha$ class at all heights.
The $\beta$ class similarly requires $v_2(v)=1$ and
$c-\beta\in16\mathbf Z_2$; at 23 one has $c-\beta\in23\mathbf Z_{23}$.

For the last class retain signed sums of the three singleton centers.
If $t$ is the square of such a sum, their quartic is
\[
 q_2(c,t,1)^2=64t e_3(c),\quad
 q_2=t^2-\tfrac{10c+50}{3}t+5c^2-30c+19.
\]
Write it as $q_2^2=q_1q_3$, with
$q_1=t(c-\theta)$ and $q_3=64e_3(c)/(c-\theta)$.
The double cover adjoining the square root of
$q_1/(\beta-\theta)$ is unramified, and its Prym is the Jacobian
of
\[
 Y^2=-(\beta-\theta)\det(M_1+2xM_2+x^2M_3),
\]
where $M_i$ are the symmetric matrices of the homogeneous quadratics.
This exact determinant specifies the twist as well as the curve.
The quartic-root field has absolute degree twelve, signature $(2,5)$
and absolute discriminant
$172707117717588564754287819295488000000$.
The explicit formula with support $2^{31}$ proves the generating
bound $48000001$, with margin above $0.080$; exact ideal relations
reduce it to 1000. The ordinary class group is $(\mathbf Z/2)^2$.
The Prym two-descent has Selmer dimension six and two-torsion
dimension one, so its rank is at most five.

A valid split triple supplies the ordered signed sums
$-a-b+d,a-b+d,-a+b+d,a+b+d$. At $c=\beta$ these are
$1/2,3/2,7/2,9/2$. Three fixed collision classes and three moving
classes have independent logarithms. For equal signs at the three
places above 23, the determinant after removing the three displacement
factors has leading term $7(c-\beta)^2$ modulo 23. The full coefficient
bounds prevent cancellation throughout the nonzero disk. For mixed
signs the simultaneous reductions at the places above 1181 give
empty intersections for the remaining three relation tests. This
forces six independent global classes, contradicting rank at most five.
The anchor $c=\beta$ itself is a collision.
The complete data are \cert{check_type13_product_covers.py},
\cert{check_type13_product_covers.m}, \cert{check_type13_alpha_selmer.m},
\cert{check_type13_beta_tail.m}, \cert{check_type13_beta_prym_rank.m}
and \cert{check_type13_beta_prym.m}. The unconditional ideal bound is
reproduced by \cert{type13_prym_class/prime_counts.cpp},
\cert{type13_prym_class/check_bound.py} and
\cert{type13_prym_class/check_factor_base.gp}.

\subsection{Type 14}\label{app:type14}
\noindent\emph{Arithmetic details in Section~\ref{detail:type14}.}\enspace
Here $(m,\mathcal A,\mathcal B)=(3,(2),(1,1,1,2))$ and
$a=(5b+3)/3$. With $K$ defined at the beginning of the appendix,
write its discriminant as $-4b h_3(b)/243$. The necessary smooth
curve is $C:y^2=-3h_3(b)$, of genus two. Its points
$\infty,(9,\pm3240)$ occupy all three residue disks at seven.
The Jacobian has rank one, certified by two-descent with Minkowski
bounds, and $D=[(9,3240)-\infty]$ is non-torsion.
On the monic model $X=-2700b$, $Y=2700^2y$, exact Cantor
arithmetic for $25D$ gives
\[
 U\equiv(X-5)^2,\quad v_7(\operatorname{disc}U)=2,
 \quad \operatorname{disc}(U)/49\equiv5,\quad7V\equiv5X+3\pmod7.
\]
The nonsquare discriminant unit resolves the two support points in
the unramified quadratic extension. Their displacement integrals
produce the rank-one logarithm and an annihilating differential
with unit coefficient in each of the three rational disks.
The linear term dominates throughout $7\mathbf Z_7$, so each disk
contains only its anchor. The finite parameter is $b=9$, and then
\[
 a=16,\qquad K(X)=(4X-1)(4X-9)(4X-81)/64.
\]
Both inputs are the same ten half-integers. The complete rank and
point calculations are \cert{check_type14_exclusion.m} and
\cert{check_type14_seven_adic.py}.

\subsection{Type 15}\label{app:type15}
\noindent\emph{Arithmetic details in Section~\ref{detail:type15}.}\enspace
The type is $(3,(1,1),(1,1,3))$. Let $c$ be the $B$ tripleton
square and $S,T$, $U,V$ the symmetric functions of the $A$ and
$B$ singleton squares. After removal of $c=1/4,9/4,25/4$, all
of which give collisions, the exact branch is
\[
 S=(20c+15)/8,\quad T=(160c^2-20c+49)/96,
 \quad U=(69-4c)/8,\quad V=(2c-23)(4c-31)/48.
\]
In addition to $c$, the following must be rational squares:
\[
\begin{aligned}
 q_T&=6(160c^2-20c+49),&q_A&=3(-80c^2+1960c+283),\\
 q_B&=3(-80c^2+808c+2875),&q_V&=3(2c-23)(4c-31).
\end{aligned}
\]
Put $K=\mathbf Q(\sqrt{15})$, $r=(245+64\sqrt{15})/20$ and
$R=(c-r)/(1/4-r)$. Its norm is a rational square. The valuations
outside 2,3,5 are even, and it is totally positive. The fundamental
unit $4+\sqrt{15}$ and the local square test at 5 show that $R$
has class 1 or 6. The identity
$(3+\sqrt{15})^2=6(4+\sqrt{15})$ identifies the second class with
the representative $c_0=49/4$. Thus only $c_0=1/4,49/4$ occur.

Use $z^2=(c_0-r)(c-r)q(c)$, taking $q=q_B$ for $1/4$ and
$q=q_V$ for $49/4$. Both cubics are rank-one elliptic curves over
$K$, with torsion orders two and four respectively. The seed point,
after two-saturation and addition of torsion, gives finite odd index.
Elliptic Chabauty returns index obstruction 4 in both cases and the
complete rational images
\[
 \{1/4,\infty\},\qquad\{31/4,23/2,49/4,\infty\}.
\]
The two extra finite values are nonsquares. At $1/4$ an $A$
singleton meets its central chain; at $49/4$ both inputs coincide.
The normalization is elementary in $K$ and every auxiliary class
proof uses a Minkowski bound. See
\cert{check_type15_unit_normalization.py} and
\cert{check_type15_exclusion.m}.

\subsection{Type 16}\label{app:type16}
\noindent\emph{Arithmetic details in Section~\ref{detail:type16}.}\enspace
The type is $(3,(1,1),(1,2,2))$. For its $B$ singleton square $c$,
the unique noncollision branch is
\[
 S=(195-20c)/12,\quad T=(160c^2-1335c+2304)/36,
 \quad U=4(9-c)/3,\quad V=2(c-6)(4c-9)/3.
\]
The vanishing elimination denominator gives only $c=1/4,25/4$
or $21/4$. The first two collide and the third is nonsquare.
The additional branches $S=c+j/4$, $j=1,9,49$, give a singleton
square equal to $c$ and hence cross-color repetitions.
Retaining $c=v^2$ and $V=y^2$ gives
\[
 y^2=\tfrac23(v^2-6)(4v^2-9).
\]
At $t=v-3/2$ the substitution $X=-10-30/t$, $Y=-30y/t^2$
gives $E:Y^2=X^3+44X^2+100X$; the exceptional branch point
$(3/2,0)$ is the identity. Its dual is $Y^2=X^3-88X^2+1536X$.
Primitive isogeny charts modulo 9,16 and 5 give descent images
$\{1,10\}$ and $\{1,6\}$, so the rank is zero. The good
reduction orders 4 and 12 at 7 and 11 bound rational torsion by
four. The points $O,(0,0),(10,\pm80)$ exhaust it. Consequently
the quartic has exactly $(v,y)=(\pm3/2,0),(0,\pm6)$ and no
rational infinity, since $8/3$ is nonsquare. At $c=0$ the
$A$ singleton discriminant is $129/16$, a nonsquare. At $c=9/4$
an $A$ singleton square is $1/4$, already in the central chain.
Every local chart, map, group count and collision is checked by
\cert{check_type16_exclusion.py}.

\subsection{Type 17}\label{app:type17}
\noindent\emph{Arithmetic details in Section~\ref{app:type17-details}.}\enspace
Here $(m,\mathcal A,\mathcal B)=(4,(1),(1,1,1,2))$ and
$4u^2=5b+1$. Put $x=1/b>0$. Both $x$ and $x+5$ are squares,
and the singleton product makes
$11025x^3-4131x^2+472x-16$ square. Thus
\[
 E_0:Y^2=X^3-4131X^2+11025\cdot472X-16\cdot11025^2,
 \qquad X=11025x
\]
contains the product point. Its rational group is infinite cyclic,
generated by $G=(11025/16,165375/64)$.
For the cubic field $K_0=\mathbf Q(w)$,
$w^3+w^2-240w+1120=0$, a root of the product polynomial is
$r=(2w^2+5w+140)/3675$. On odd multiples $x-r$ has class
$d=1/16-r$. The curve $Y^2=X(X+5d)(X-rd)$ has rank one,
full two-torsion, and an explicitly two-saturated non-torsion point.
Elliptic Chabauty returns index obstruction 16 and images
$\infty,0,-5,1/16$. Positivity and finiteness leave $b=16$,
whose doubleton meets the central octupleton. This handles all
odd multiples.

For an even multiple write the product point as $2Q$, set $T=X(Q)$,
$a=11025$ and
\[
\begin{aligned}
 g(T)&=T^3-4131T^2+472aT-16a^2,\\
 A(T)&=T^4-944aT^2+128a^2T+a^2(472^2-4\cdot4131\cdot16),\\
 B(T)&=A(T)+20ag(T).
\end{aligned}
\]
Duplication gives $x=A/(4ag)$ and $x+5=B/(4ag)$; all three
of $g,A,B$ are nonzero squares. The two-cover descents on
$y^2=gA$ and $y^2=gB$ have sizes two and one. Compatibility
in their common cubic factor selects the identity in both, so
$T-\theta$ is square for a root $\theta$ of $A$, and
$Q\in2E_0(\mathbf Q)$. Its formal parameter $s=-T/Y(Q)$ then
satisfies $k=v_2(s)\geq3$ and $v_2(T)=-2k$.

Write $K=\mathbf Q(w)$,
$w^4-194w^2-96w+8302=0$, and
$\theta=3(2w^3+4w^2-240w-532)$. The monic elliptic curve
$Y'^2=h(T)=A(T)/(T-\theta)$ contains the putative point.
Its two-division root field has degree twelve, is totally real,
and has discriminant $58160039859394517078535559761100800000$.
The explicit formula at support $2^{20}$ has margin greater than
$0.262$. Exact ideal reduction to norm 1000 certifies class group
$\mathbf Z/2$ and Selmer dimension two.
At the unique totally ramified quartic dyadic place, put
$m=105\cdot2^{12}$. Complete local Kummer images show that
$m\log E(K_v)=2^{11}\mathfrak p$. In its basis
$(2^{12},2^{11}\pi,2^{11}\pi^2,2^{11}\pi^3)$ the global Selmer
image has dimension two and does not contain $(1,0,0,0)$.
The ratio of the two forced formal parameters satisfies
$(z/s)^2=g(T)/h(T)\in1+2^{2k}\mathcal O_{K_v}$.
The normalized global logarithm therefore has precisely that excluded
primitive vector. Dividing by the largest global power of two,
using injectivity of the localized Selmer image, proves the
contradiction at every height.
The odd calculation is \cert{check_type17_odd_branch.m}; the
unconditional field bound and the complete dyadic argument are
\cert{check_type17_factor_base.gp} and
\cert{check_type17_even_branch.m}, together with the
\texttt{--type17} option of \cert{type13_prym_class/check_bound.py}.
